The coding-agent study completes 80 attempts: ten selected Astropy tasks, twice each under \sys{}, native MTP-5, SGLang EAGLE, and plain autoregressive decoding of the same checkpoint on the same vLLM build as a control. All arms use Qwen Code 0.19.4, the same checkpoint, a 24,000-token response cap, and a 9,000-second task budget. The agent inspects repositories, edits files, and executes tools. Table~\ref{tab:agent-study} summarizes all attempts, including seven empty patches; Appendix~\ref{app:workload} gives task-level results. All arms use presence penalty 1.0; the \sys{} arm applies per-path penalty histories (Section~\ref{sec:sampling-validation}).

For completed request $r$, let $n_r$ be the output-token count, $\ell_r$ the server-recorded end-to-end latency, and $f_r$ the time to first output. We report the pooled decode rate
\begin{equation}
 D_{\mathrm{pool}} =
 \frac{\sum_{r=1}^{R}(n_r-1)}{\sum_{r=1}^{R}(\ell_r-f_r)}.
 \label{eq:pooled-decode}
\end{equation}
Counts and durations cover the same completed requests, including requests that summarize an overlong agent context. The numerator counts output intervals; the denominator excludes first-output latency and between-request agent/tool work. Thus $D_{\mathrm{pool}}$ measures token production over accumulated request durations, rather than aggregate wall-clock throughput.

\begin{table}[t]
\caption{Ten-task coding-agent study, two attempts per task and arm. Resolved counts are per attempt (of ten); requests, tokens, and minutes sum both attempts; decode rates pool both attempts with server-counter differences in Eq.~\ref{eq:pooled-decode}. Plain decoding is the same-build control without speculation. Minutes sum agent durations, including tools but excluding server startup and evaluation.}
\label{tab:agent-study}
\centering\small
\setlength{\tabcolsep}{3pt}
\begin{tabular}{@{}lcrrrr@{}}
\toprule
Arm & Resolved & Requests & Tokens & Tokens/s & Minutes \\
\midrule
\sys{} & 5, 6 & 457 & 440,304 & 26.62 & 323.3 \\
Native MTP-5 & 6, 6 & 463 & 302,942 & 24.95 & 240.7 \\
SGLang EAGLE & 5, 5 & 475 & 386,367 & 28.46 & 261.0 \\
Plain decoding & 7, 7 & 481 & 271,910 & 10.47 & 464.7 \\
\bottomrule
\end{tabular}
\end{table}

\sys{} resolves five then six tasks, native MTP-5 six and six, SGLang five and five, and plain decoding seven and seven; empty patches (three for \sys{}, two each for MTP-5 and SGLang, none for plain decoding, over both attempts) count as failures without invoking the test evaluator. Over twenty attempts per arm, plain decoding resolves 14, MTP-5 12, \sys{} 11, and SGLang 10; a two-sided Fisher exact test on 14 against 10 of 20 gives $p=0.33$, so the study does not separate the arms on outcome. These are descriptive outcomes on the selected tasks, not benchmark-wide quality estimates.

The live agent trajectories diverge after the first tool call, and output volume varies as much between attempts of one arm as between arms: \sys{} emits 258,889 then 181,415 output tokens, MTP-5 134,643 then 168,299, SGLang 158,827 then 227,540, and plain decoding 98,477 then 173,433. The first attempt's ordering, in which \sys{} emitted the most tokens and ran longest, does not repeat: in the second attempt \sys{} emits 8\% more tokens than MTP-5 and finishes in fewer minutes. Pooled over both attempts, \sys{} decodes at 26.62 tokens/s against 24.95 for MTP-5, 28.46 for SGLang, and 10.47 for plain decoding; MTP-5's second attempt accepted 2.84 drafts per event, below the 3.17 of its fixed-request replays, which lowers its live rate. Request counts, prompt lengths, and tool calls differ across arms and attempts, so sharing tasks and agent settings does not isolate a verifier's effect on task completion.
